\documentclass[11pt,a4paper]{article}

\usepackage[margin=17mm]{geometry}
\usepackage{fontspec}
\usepackage{enumitem}
\usepackage{hyperref}
\usepackage{titlesec}
\usepackage{microtype}
\usepackage{needspace}

\setmainfont{Latin Modern Roman}
\hypersetup{colorlinks=true,urlcolor=black}
\setlength{\parindent}{0pt}
\setlength{\parskip}{2pt}
\pagestyle{empty}
\setlist[itemize]{leftmargin=*,nosep,topsep=2pt}
\titleformat{\section}{\large\bfseries}{}{0pt}{}[\titlerule]
\titlespacing*{\section}{0pt}{9pt}{4pt}

\newcommand{\role}[3]{%
  \Needspace{6\baselineskip}%
  \textbf{#1}\hfill #3\\
  \textit{#2}\par
}

\begin{document}

\noindent
\begin{minipage}[c]{0.52\textwidth}
  {\fontsize{24}{27}\selectfont\bfseries Ajinkya Kamat}\par

\end{minipage}%
\begin{minipage}[c]{0.48\textwidth}
  \raggedleft\small
  Mumbai, India\par
  \href{mailto:kamat.ajinkya@gmail.com}{kamat.ajinkya@gmail.com}\par
  \href{https://ajinkyakamat.github.io}{ajinkyakamat.github.io}
\end{minipage}

\section*{Profile}
Electrical engineer and roboticist with 10+ years of experience delivering
robotic systems from electronics and motor drives through embedded firmware,
real-time control, integration and field deployment. Technical lead in
cross-disciplinary engineering, field reliability and research-to-system integration.

\section*{Professional Experience}

\role{FieldAI}{Electrical and Firmware Lead}{Feb 2024 -- Present}
\begin{itemize}
  \item Lead electrical and firmware development of Falcon and Ferrite
  autonomy payloads, integrating LiDAR, depth sensors, onboard compute
  and power systems for mobile-robot deployment.
  \item Helped scale deployed robotic systems from approximately 3 to
  100+ through payload redesign, system integration and validation.
  \item Extended reported payload mean time between failures from about
  3 days to 2 months through field-failure analysis and electronics
  and firmware redesign.
  \item Reduced payload iteration cycles from approximately 3 months to
  1 month by introducing structured engineering workflows.
  \item Coordinate design, validation and deployment across distributed
  engineering teams in India and the United States.
\end{itemize}

\role{Carnegie Mellon University, Biorobotics Lab}{Graduate Student Researcher}{Oct 2021 -- Dec 2023}
\begin{itemize}
  \item Integrated ROS and MoveIt across three lab robots by developing
  control nodes, drivers, robot descriptions and motion-planning packages;
  maintained Docker environments and MuJoCo simulation for reproducible
  research across the lab.
  \item Implemented distributed CANopen actuator control, enabling
  coordinated motion on a novel 6-DOF direct-drive Delta manipulator.
  \item Developed a custom direct-drive gripper with a directionally compliant
  controller, enabling grasps of low-profile objects such as flat-lying phones.
  \item Applied Demons deformable image registration to estimate screw
  locations after smartphone deformation for robotic disassembly.
\end{itemize}

\role{Systemantics India Pvt. Ltd.}{Robotics Systems Engineering}{2020 -- 2021}
\begin{itemize}
  \item Modernized the industrial robot system controller by evaluating
  ARMv8 hardware and developing a modular ROS 2 software architecture.
  \item Achieved reported sub-10\,$\mu$s control latency by optimizing
  PREEMPT\_RT Linux; established Git, Docker and apt-based deployment workflows.
\end{itemize}

\role{University of Washington, Personal Robotics Lab}{Visiting Scientist, RCTA Program}{Apr 2018 -- May 2020}
\begin{itemize}
  \item Delivered the capstone unstructured manipulation demonstration
  for RoMan's decade-long, seven-institution research program by
  integrating its research software packages and developing hand--eye
  calibration for perception-guided debris clearing.
  \item Built experimental data-collection infrastructure for a
  chopstick telemanipulation study, supporting research published at IROS 2020.
\end{itemize}

\newpage

\role{Systemantics India Pvt. Ltd.}{Robotics Electronics and Firmware Engineering}{2014 -- 2018}
\begin{itemize}
  \item Led a three-engineer electronics team in end-to-end development
  of controllers and firmware for parallel and articulated industrial robots.
  \item Reduced the cost of a 1.2\,kW three-phase brushless motor driver
  by approximately 30\% through electrical redesign; developed
  FreeRTOS-based dual-motor-controller firmware.
  \item Developed DSP-based joint-space and task-space motion control
  for coordinated operation of industrial manipulators.
  \item Reduced harmonic-drive-induced robot vibrations by approximately
  50\% through transmission-error diagnosis and control improvements.
\end{itemize}

\role{AVL}{Engineering Intern}{May 2013 -- Jul 2013}
\begin{itemize}
  \item Replaced hardware PLC-based closed-loop temperature, fluid and
  coolant control on engine test beds with NI LabVIEW software on Windows CE;
  AVL subsequently offered the solution as a commercial add-on.
\end{itemize}

\section*{Leadership — IvLabs, VNIT}
\textit{Co-founder and Technical Mentor \hfill 2013 -- Present}
\begin{itemize}
  \item Co-founded a student-led robotics research lab that has grown
  to 80+ members, with 35+ research papers and 5+ patent filings produced
  by successive cohorts.
  \item Originated research directions, formed project teams and mentored
  students through Swayat, Sahayak, ReQuBiS, OSMOS and prosthetic-hand projects.
  \item Led ReBiS research; established the Swayat team, built its first
  biped prototype and mentored the second-generation humanoid.
\end{itemize}

\section*{Technical Skills}
\textbf{Electronics and controls:} Circuit and PCB design, power electronics,
motor drives, control systems, hardware bring-up, instrumentation,
failure analysis\\
\textbf{Embedded and robotics:} C, C++, Python, FreeRTOS, CANopen,
PREEMPT\_RT Linux, ROS, ROS 2, MoveIt, motor and motion control\\
\textbf{Engineering tools:} KiCad, OrCAD, SPICE, MATLAB, Git, Docker,
NI LabVIEW, Linux

\section*{Selected Research Publications}
Co-author: ReBiS (IROS 2014); \textit{Telemanipulation with Chopsticks}
(IROS 2020); \textit{Toward Fieldable Human-Scale Mobile Manipulation
Using RoMan} (SPIE 2020).

\section*{Education}
\textbf{Visvesvaraya National Institute of Technology (VNIT), Nagpur}\\
B.Tech., Electronics and Communication Engineering, May 2014

\end{document}
